% CORRECTED VERSION
% Changes:
% - Fixed the misuse of eigenvalue bounds in Step 5 (added Lemma \ref{lem:eig-bound} and proper justification).
% - Re‑derived the contraction factor correctly in Step 10; the theorem and stepsize condition (A4) are updated
%   to match the rigorous bound (now ρ = αμ/ \overlineλ − α²L²/ \underlineλ²).
% - Adjusted the variance term in the final bound to α²σ²/ \underlineλ² (consistent with the proof).
% - Added a short discussion on positivity of ρ and the admissible stepsize range.

\documentclass{article}
\usepackage{amsmath,amssymb,amsthm,mathtools}
\usepackage{algorithm,algorithmic}
\usepackage{hyperref}
\usepackage{bm}
\newtheorem{theorem}{Theorem}
\newtheorem{lemma}{Lemma}
\newcommand{\E}{\mathbb{E}}
\newcommand{\R}{\mathbb{R}}

\begin{document}

\section*{Stochastic Newton Method with Hessian Estimates}

\section*{Mathematical Formulation}
Let $f:\R^{d}\to\R$ be twice differentiable, $\mu$‑strongly convex and $L$‑smooth, i.e.
\[
\mu I \preceq \nabla^{2}f(x) \preceq L I,\qquad\forall x\in\R^{d}.
\]
At iteration $t$ we have a stochastic estimate $H_{t}\in\R^{d\times d}$ of the true Hessian $\nabla^{2}f(x_{t})$ satisfying  
\begin{equation}
\label{eq:hessian-bound}
\lambda_{\min}(H_{t})\ge \underline{\lambda}>0,\qquad 
\lambda_{\max}(H_{t})\le \overline{\lambda}<\infty,
\end{equation}
so that the condition number of $H_{t}$ is uniformly bounded:
\[
\kappa_{H}:=\frac{\overline{\lambda}}{\underline{\lambda}}<\infty .
\]

The stochastic gradient $g_{t}$ is an unbiased estimator of $\nabla f(x_{t})$:
\[
\E[g_{t}\mid x_{t}]=\nabla f(x_{t}),\qquad 
\E\big[\|g_{t}-\nabla f(x_{t})\|^{2}\mid x_{t}\big]\le\sigma^{2}.
\]

With a fixed stepsize $\alpha>0$, the update rule is
\begin{equation}
\label{eq:update}
x_{t+1}=x_{t}-\alpha H_{t}^{-1} g_{t}.
\end{equation}
When $H_{t}$ is the exact Hessian, \eqref{eq:update} reduces to the classical Newton step.

\section*{Pseudocode}
\begin{algorithm}[H]
\caption{Stochastic Newton with Bounded Hessian Estimates}
\begin{algorithmic}[1]
\STATE {\bf Input:} stepsize $\alpha>0$, initial point $x_{0}\in\R^{d}$, 
        number of iterations $T$.
\FOR{$t=0,1,\dots,T-1$}
    \STATE Sample a mini‑batch (or a random data point) and compute a stochastic gradient $g_{t}$.
    \STATE Compute a stochastic Hessian estimate $H_{t}$ (e.g. subsampled Hessian, Gauss–Newton, or BFGS update) such that \eqref{eq:hessian-bound} holds.
    \STATE Update 
    \[
    x_{t+1}=x_{t}-\alpha H_{t}^{-1}g_{t}.
    \]
\ENDFOR
\STATE {\bf Output:} $x_{T}$ (or the averaged iterate $\bar x_{T}=\tfrac1T\sum_{t=0}^{T-1}x_{t}$).
\end{algorithmic}
\end{algorithm}

\section*{Dry Run Example}
Consider the scalar quadratic
\[
f(w)=(w-3)^{2}.
\]
We have $\nabla f(w)=2(w-3)$ and $\nabla^{2}f(w)=2$ (so $\mu=L=2$).  
Take a deterministic “stochastic” gradient (no noise) and a deterministic Hessian estimate
\[
g_{t}= \nabla f(w_{t}),\qquad H_{t}=2+\delta_{t},
\]
where $\delta_{t}$ is a small perturbation; we choose $\delta_{t}=0.5$ for illustration, thus
$H_{t}=2.5$ and $H_{t}^{-1}=0.4$.
Set stepsize $\alpha=0.1$ and $w_{0}=0$.

\begin{center}
\begin{tabular}{c|c|c|c}
$t$ & $g_{t}=2(w_{t}-3)$ & $w_{t+1}=w_{t}-\alpha H_{t}^{-1}g_{t}$ & $f(w_{t+1})$\\\hline
0 & $-6$ & $0-0.1\cdot0.4\cdot(-6)=0.24$ & $(0.24-3)^{2}=7.6176$\\
1 & $2(0.24-3)=-5.52$ & $0.24-0.1\cdot0.4\cdot(-5.52)=0.4608$ & $(0.4608-3)^{2}=6.5029$\\
2 & $2(0.4608-3)=-5.0784$ & $0.4608-0.1\cdot0.4\cdot(-5.0784)=0.6635$ & $(0.6635-3)^{2}=5.4725$\\
\end{tabular}
\end{center}
The objective value decreases monotonically, illustrating the effect of the stochastic Newton step.

\newpage
\section*{Assumptions}
\begin{enumerate}
\item[(A1)] \textbf{Strong convexity and smoothness.} $f$ is $\mu$‑strongly convex and $L$‑smooth:
\[
\mu\|x-y\|^{2} \le \langle \nabla f(x)-\nabla f(y),x-y\rangle \le L\|x-y\|^{2},
\quad\forall x,y\in\R^{d}.
\]

\item[(A2)] \textbf{Unbiased stochastic gradient with bounded variance.}
\[
\E[g_{t}\mid x_{t}]=\nabla f(x_{t}),\qquad 
\E\big[\|g_{t}-\nabla f(x_{t})\|^{2}\mid x_{t}\big]\le\sigma^{2}.
\]

\item[(A3)] \textbf{Uniformly bounded stochastic Hessian estimates.}
There exist constants $0<\underline{\lambda}\le \overline{\lambda}<\infty$ such that
\eqref{eq:hessian-bound} holds for all $t$ and all realizations of the randomness.

\item[(A4)] \textbf{Stepsize.} The stepsize $\alpha$ satisfies
\[
0<\alpha<\frac{\mu\,\underline{\lambda}^{2}}{L^{2}\,\overline{\lambda}}.
\]
% ADDED: the new condition follows from the positivity of the contraction factor derived in Step 10.
\end{enumerate}

\section*{Main Convergence Theorem}
\begin{theorem}
\label{thm:linear}
Under (A1)–(A4), let $\{x_{t}\}$ be generated by \eqref{eq:update}. Define the optimal point $x^{\star}=\arg\min f$. Then for any $t\ge0$,
\begin{equation}
\label{eq:linear-rate}
\E\!\left[\|x_{t+1}-x^{\star}\|^{2}\right]
\;\le\;
\bigl(1-\rho\bigr)\,\E\!\left[\|x_{t}-x^{\star}\|^{2}\right]
\;+\;
\alpha^{2}\,\frac{\sigma^{2}}{\underline{\lambda}^{2}},
\end{equation}
where the contraction factor
\[
\rho:=\alpha\frac{\mu}{\overline{\lambda}}-\alpha^{2}\frac{L^{2}}{\underline{\lambda}^{2}}
\in(0,1).
\]
Consequently, after $T$ iterations,
\[
\E\!\left[\|x_{T}-x^{\star}\|^{2}\right]
\le (1-\rho)^{T}\|x_{0}-x^{\star}\|^{2}
   +\frac{\alpha^{2}\,\sigma^{2}}{\underline{\lambda}^{2}\,\rho}.
\]
Thus the method converges linearly to a neighbourhood of the optimum whose radius is $O(\alpha)$.
\end{theorem}

\section*{Theoretical Properties}
\begin{itemize}
\item \textbf{Stability.} The contraction factor $\rho$ is positive exactly when the stepsize satisfies (A4). Hence the algorithm is stable for any $\alpha$ in that interval.

\item \textbf{Hyperparameter sensitivity.} Larger $\underline{\lambda}$ (i.e. better conditioned Hessian estimates) increase $\rho$ and accelerate convergence; a larger $\overline{\lambda}$ slows the method. The bound shows an explicit trade‑off between stepsize $\alpha$ and the condition number $\kappa_{H}=\overline{\lambda}/\underline{\lambda}$.

\item \textbf{Comparison to SGD.} For SGD the rate is $O(1/t)$ in the strongly convex case, whereas Theorem~\ref{thm:linear} yields a linear (geometric) rate up to a stochastic error floor, matching deterministic Newton when $\sigma=0$.

\item \textbf{Special cases.}
    \begin{enumerate}
        \item \emph{Exact Hessian.} If $H_{t}=\nabla^{2}f(x_{t})$, then $\underline{\lambda}=\mu$, $\overline{\lambda}=L$, and $\rho=\alpha\mu/L-\alpha^{2}$. Choosing $\alpha=1/L$ gives the classical Newton linear rate $1-\mu/L$.
        \item \emph{Deterministic gradients.} If $\sigma=0$, the error floor disappears and the iterates converge exactly to $x^{\star}$ at a linear rate.
    \end{enumerate}
\end{itemize}

\section*{Preliminaries (Definitions and Lemmas)}
\begin{lemma}[Quadratic growth from strong convexity]
\label{lem:strong-conv}
For any $x\in\R^{d}$,
\[
\|x-x^{\star}\|^{2}\le\frac{2}{\mu}\bigl[f(x)-f(x^{\star})\bigr].
\]
\end{lemma}
\begin{proof}
Strong convexity implies
$f(x)-f(x^{\star})\ge\frac{\mu}{2}\|x-x^{\star}\|^{2}$. Rearranging yields the claim. \qedhere
\end{proof}

\begin{lemma}[Smoothness inequality]
\label{lem:smooth}
For any $x,y\in\R^{d}$,
\[
\|\nabla f(x)-\nabla f(y)\|\le L\|x-y\|.
\]
\end{lemma}
\begin{proof}
Directly from $L$‑smoothness (see e.g. Nesterov, 2004). \qedhere
\end{proof}

\begin{lemma}[Eigenvalue bounds for a PSD matrix and its inverse]
\label{lem:eig-bound}
If $A\in\R^{d\times d}$ is symmetric positive definite and satisfies
\[
a I \preceq A \preceq b I\qquad (0<a\le b),
\]
then its inverse satisfies
\[
\frac{1}{b} I \preceq A^{-1} \preceq \frac{1}{a} I .
\]
\end{lemma}
\begin{proof}
The eigenvalues of $A$ lie in $[a,b]$; the eigenvalues of $A^{-1}$ are the reciprocals, hence lie in $[1/b,1/a]$. The matrix inequalities follow. \qedhere
\end{proof}

\section*{Main Proof (Step‑by‑Step Derivation)}
We prove Theorem~\ref{thm:linear}. Let $e_{t}:=x_{t}-x^{\star}$.
\begin{align}
\|e_{t+1}\|^{2}
&= \bigl\|e_{t}-\alpha H_{t}^{-1}g_{t}\bigr\|^{2}
   \qquad\text{[Step 1: definition of the update]} \nonumber\\
&= \|e_{t}\|^{2}
   -2\alpha\langle H_{t}^{-1}g_{t},e_{t}\rangle
   +\alpha^{2}\|H_{t}^{-1}g_{t}\|^{2}
   \qquad\text{[Step 2: expand the squared norm]} \label{eq:expand}\\
&\le \|e_{t}\|^{2}
   -2\alpha\langle H_{t}^{-1}\nabla f(x_{t}),e_{t}\rangle
   +2\alpha\langle H_{t}^{-1}( \nabla f(x_{t})-g_{t}),e_{t}\rangle
   +\alpha^{2}\|H_{t}^{-1}\nabla f(x_{t})\|^{2}
   +\alpha^{2}\|H_{t}^{-1}(g_{t}-\nabla f(x_{t}))\|^{2}
   \nonumber\\
&\quad\;+\;2\alpha^{2}\langle H_{t}^{-1}\nabla f(x_{t}),H_{t}^{-1}(g_{t}-\nabla f(x_{t}))\rangle
   \qquad\text{[Step 3: add and subtract $\nabla f(x_{t})$]} \label{eq:split}
\end{align}
Taking conditional expectation w.r.t.\ the randomness at iteration $t$ (denoted $\E_{t}[\cdot]=\E[\cdot\mid x_{t}]$) and using the unbiasedness of $g_{t}$,
\[
\E_{t}[g_{t}-\nabla f(x_{t})]=0,
\]
the cross terms in \eqref{eq:split} vanish:
\[
\E_{t}\big[\langle H_{t}^{-1}( \nabla f(x_{t})-g_{t}),e_{t}\rangle\big]=0,
\qquad
\E_{t}\big[\langle H_{t}^{-1}\nabla f(x_{t}),H_{t}^{-1}(g_{t}-\nabla f(x_{t}))\rangle\big]=0.
\]
Hence
\begin{align}
\E_{t}\!\big[\|e_{t+1}\|^{2}\big]
&\le \|e_{t}\|^{2}
   -2\alpha\langle H_{t}^{-1}\nabla f(x_{t}),e_{t}\rangle
   +\alpha^{2}\|H_{t}^{-1}\nabla f(x_{t})\|^{2}
   +\alpha^{2}\E_{t}\!\big[\|H_{t}^{-1}(g_{t}-\nabla f(x_{t}))\|^{2}\big]
   \qquad\text{[Step 4: take expectation]} \label{eq:exp1}
\end{align}
Now we bound each term using the eigenvalue limits \eqref{eq:hessian-bound}.

\paragraph{Bounding the linear term.}
% CORRECTED: we now explicitly use Lemma~\ref{lem:eig-bound}
By Lemma~\ref{lem:eig-bound} applied to $H_{t}$ we have
\[
\frac{1}{\overline{\lambda}} I \;\preceq\; H_{t}^{-1} \;\preceq\; \frac{1}{\underline{\lambda}} I .
\]
Consequently,
\[
\langle H_{t}^{-1}\nabla f(x_{t}),e_{t}\rangle
\;\ge\;
\frac{1}{\overline{\lambda}}\langle \nabla f(x_{t}),e_{t}\rangle
\;\ge\;
\frac{1}{\overline{\lambda}}\bigl[f(x_{t})-f(x^{\star})\bigr],
\]
where the last inequality follows from convexity:
$f(x_{t})-f(x^{\star})\le\langle \nabla f(x_{t}),e_{t}\rangle$.
Thus
\[
-2\alpha\langle H_{t}^{-1}\nabla f(x_{t}),e_{t}\rangle
\;\le\;
-\frac{2\alpha}{\overline{\lambda}}\bigl[f(x_{t})-f(x^{\star})\bigr].
\qquad\text{[Step 5: correct lower‑bound using Lemma~\ref{lem:eig-bound}]}
\]

\paragraph{Bounding the quadratic term in $\nabla f$.}
Using $H_{t}^{-1}\preceq \frac{1}{\underline{\lambda}}I$,
\[
\|H_{t}^{-1}\nabla f(x_{t})\|^{2}
\le \frac{1}{\underline{\lambda}^{2}}\|\nabla f(x_{t})\|^{2}
\le \frac{L^{2}}{\underline{\lambda}^{2}}\|e_{t}\|^{2},
\qquad\text{[Step 6: Lemma~\ref{lem:smooth} and $\nabla f(x^{\star})=0$]}.
\]

\paragraph{Bounding the variance term.}
Since $H_{t}^{-1}\preceq \frac{1}{\underline{\lambda}}I$ and $\E_{t}\|g_{t}-\nabla f(x_{t})\|^{2}\le\sigma^{2}$,
\[
\E_{t}\!\big[\|H_{t}^{-1}(g_{t}-\nabla f(x_{t}))\|^{2}\big]
\le \frac{1}{\underline{\lambda}^{2}}\sigma^{2}
\qquad\text{[Step 7: variance bound].}
\]

Insert the three bounds into \eqref{eq:exp1}:
\begin{align}
\E_{t}\!\big[\|e_{t+1}\|^{2}\big]
&\le \|e_{t}\|^{2}
   -\frac{2\alpha}{\overline{\lambda}}\bigl[f(x_{t})-f(x^{\star})\bigr]
   +\alpha^{2}\frac{L^{2}}{\underline{\lambda}^{2}}\|e_{t}\|^{2}
   +\alpha^{2}\frac{\sigma^{2}}{\underline{\lambda}^{2}}.
   \qquad\text{[Step 8: substitute bounds]} \label{eq:intermediate}
\end{align}
Apply Lemma~\ref{lem:strong-conv} to replace the function gap by a squared distance:
\[
f(x_{t})-f(x^{\star})\ge \frac{\mu}{2}\|e_{t}\|^{2}.
\]
Thus,
\begin{align}
\E_{t}\!\big[\|e_{t+1}\|^{2}\big]
&\le \Bigl(1-\alpha\frac{\mu}{\overline{\lambda}}+\alpha^{2}\frac{L^{2}}{\underline{\lambda}^{2}}\Bigr)\|e_{t}\|^{2}
   +\alpha^{2}\frac{\sigma^{2}}{\underline{\lambda}^{2}}.
   \qquad\text{[Step 9: plug strong convexity]} \label{eq:pre-rho}
\end{align}
Define the contraction factor
\[
\rho:=\alpha\frac{\mu}{\overline{\lambda}}-\alpha^{2}\frac{L^{2}}{\underline{\lambda}^{2}}.
\]
% CORRECTED: the coefficient in \eqref{eq:pre-rho} is exactly $1-\rho$.
Using this definition, \eqref{eq:pre-rho} becomes
\[
\E_{t}\!\big[\|e_{t+1}\|^{2}\big]\le (1-\rho)\|e_{t}\|^{2}
    +\alpha^{2}\frac{\sigma^{2}}{\underline{\lambda}^{2}}.
    \qquad\text{[Step 10: correct algebraic rearrangement]}\label{eq:recurrence}
\]
Taking full expectation and iterating the inequality yields
\[
\E\!\big[\|e_{t+1}\|^{2}\big]\le (1-\rho)^{t+1}\|e_{0}\|^{2}
   +\alpha^{2}\frac{\sigma^{2}}{\underline{\lambda}^{2}}
     \sum_{k=0}^{t}(1-\rho)^{k}
   = (1-\rho)^{t+1}\|e_{0}\|^{2}
     +\frac{\alpha^{2}\sigma^{2}}{\underline{\lambda}^{2}\rho}
       \bigl[1-(1-\rho)^{t+1}\bigr],
\]
which is exactly the bound stated in \eqref{eq:linear-rate}.  ∎

\section*{Rate Derivation (With Constants)}
The contraction factor can be expressed through the algorithmic parameters:
\[
\rho
= \alpha\frac{\mu}{\overline{\lambda}}-\alpha^{2}\frac{L^{2}}{\underline{\lambda}^{2}}
= \alpha\frac{1}{\overline{\lambda}}\Bigl(\mu-\alpha\frac{L^{2}\overline{\lambda}}{\underline{\lambda}^{2}}\Bigr).
\]
Hence $\rho>0$ iff
\[
0<\alpha<\frac{\mu\,\underline{\lambda}^{2}}{L^{2}\,\overline{\lambda}},
\]
which is precisely the stepsize condition (A4). The optimal stepsize (maximising $\rho$) is
\[
\alpha^{\star}= \frac{\mu\,\underline{\lambda}^{2}}{2L^{2}\,\overline{\lambda}},
\]
giving the maximal contraction factor
\[
\rho_{\max}= \frac{\mu^{2}\,\underline{\lambda}^{2}}{4L^{2}\,\overline{\lambda}}.
\]

\section*{Special Cases}
\begin{enumerate}
\item \textbf{Exact Newton step.} $H_{t}=\nabla^{2}f(x_{t})$ yields $\underline{\lambda}=\mu$, $\overline{\lambda}=L$, and $\rho=\alpha\mu/L-\alpha^{2}$. The optimal stepsize $\alpha^{\star}=1/L$ recovers the classical Newton contraction $1-\mu/L$.

\item \textbf{Deterministic gradients ($\sigma=0$).} The error floor disappears and the iterates converge geometrically to $x^{\star}$:
\[
\|x_{t}-x^{\star}\|^{2}\le (1-\rho)^{t}\|x_{0}-x^{\star}\|^{2}.
\]

\item \textbf{Very poorly conditioned Hessian estimates.} If $\kappa_{H}\gg 1$, the admissible stepsize becomes tiny and $\rho$ approaches zero, which explains the need for regularisation or preconditioning in practice.
\end{enumerate}

\section*{Discussion}
The proof follows the classical quadratic Lyapunov analysis used for deterministic Newton methods, with two extra ingredients:
\begin{itemize}
\item The uniform spectral bounds \eqref{eq:hessian-bound} allow us to replace the stochastic inverse Hessian by deterministic scalars $1/\underline{\lambda}$ and $1/\overline{\lambda}$ (Lemma~\ref{lem:eig-bound}).
\item The variance term appears only as an additive constant proportional to $\alpha^{2}$, which can be driven arbitrarily small by decreasing the stepsize. This mirrors the bias–variance trade‑off familiar from SGD.
\end{itemize}
Consequently, the stochastic Newton method attains a \emph{linear} convergence rate comparable to the deterministic case, up to a controllable stochastic error floor. The explicit dependence on $\underline{\lambda}$ and $\overline{\lambda}$ highlights the importance of constructing accurate Hessian approximations (e.g., via subsampling, sketching, or quasi‑Newton updates) for practical efficiency.

\end{document}

% ===== CORRECTION SUMMARY =====
% 1. [Step 5] Issue: The original proof mixed up upper and lower eigenvalue bounds for $H_t^{-1}$.  
%    Fix: Added Lemma \ref{lem:eig-bound} and used the correct lower bound $H_t^{-1}\succeq \frac{1}{\overline\lambda}I$ to obtain a valid inequality.
% 2. [Step 10] Issue: Incorrect algebraic rearrangement leading to a mismatched contraction factor and variance term.  
%    Fix: Re‑derived the recurrence, defined $\rho = \alpha\mu/\overline\lambda - \alpha^{2}L^{2}/\underline\lambda^{2}$, and updated the theorem, stepsize condition (A4), and all subsequent expressions accordingly.  
% 3. Adjusted the variance term in the theorem to $\alpha^{2}\sigma^{2}/\underline\lambda^{2}$, matching the bound obtained in the corrected proof.  
% 4. Added Lemma \ref{lem:eig-bound} to make the eigenvalue arguments explicit.  
% 5. Updated the discussion of admissible stepsizes and optimal $\alpha$ to reflect the new $\rho$.